\chapter{Consolidated Configuration Reference}
\label{app:configuration}

This appendix reproduces a complete, consolidated \texttt{dagon.ini}
reference, assembling every section discussed individually in
\Cref{ch:configuration}, for use as a single lookup document while authoring
workflows.

\section{How to Use This Appendix}

Two reading strategies suit two different needs. A reader setting up a new
workflow's configuration for the first time should treat the consolidated
listing below as a \emph{starting template}: copy it, delete every section
the workflow's task types do not actually require (per the
backend-selection guidance of \Cref{tab:backend-selection}), and fill in
only the remaining, genuinely needed sections --- resisting the temptation
to retain unused sections ``just in case,'' since an unused but populated
section is, per \Cref{ch:configuration}'s auditing-procedure discussion,
one more place a stray credential placeholder could later be mistaken for a
real value. A reader instead debugging an existing, unfamiliar
\texttt{dagon.ini} should use the quick-reference table at the end of this
appendix as an index, then follow each entry's chapter cross-reference back
to \Cref{ch:configuration}'s fuller discussion of that section's semantics
and failure modes.

\begin{lstlisting}[style=ini, caption={A consolidated, illustrative \texttt{dagon.ini} (placeholders only --- never commit real secrets).}]
; ---------------------------------------------------------------
; Optional DAGon service (disabled by default; see Sec. 5.2.1)
; ---------------------------------------------------------------
[dagon_service]
route = http://localhost:57000
use = False

; ---------------------------------------------------------------
; FTP publishing / transversal data download (Sec. 5.2.2)
; ---------------------------------------------------------------
[ftp_pub]
ip = localhost

; ---------------------------------------------------------------
; Service/node addressing (Sec. 5.2.3)
; ---------------------------------------------------------------
[dagon_ip]
ip = localhost

; ---------------------------------------------------------------
; Local batch execution and scratch layout (Sec. 5.2.4)
; ---------------------------------------------------------------
[batch]
scratch_dir_base=/tmp/
run_base=
remove_dir=False
threads=1

; ---------------------------------------------------------------
; Slurm scheduler defaults (Sec. 5.2.5)
; ---------------------------------------------------------------
[slurm]
partition=

; ---------------------------------------------------------------
; LLM providers: any number of [llm.<name>] sections (Sec. 5.2.6)
; ---------------------------------------------------------------
[llm.example]
endpoint=https://api.example.org
api_key_env=DAGON_LLM_EXAMPLE_API_KEY
model=example-chat-model

; ---------------------------------------------------------------
; Cloud providers (Sec. 5.2.8)
; ---------------------------------------------------------------
[ec2]
key=<aws_access_key_id>
secret=<aws_secret_access_key>
region=

[digitalocean]
key=

[gce]
key=
secret=
project=

; ---------------------------------------------------------------
; Standard Python logging configuration (Sec. 5.2.9)
; ---------------------------------------------------------------
[loggers]
keys=root

[handlers]
keys=stream_handler

[formatters]
keys=formatter

[logger_root]
level=DEBUG
handlers=stream_handler

[handler_stream_handler]
class=StreamHandler
level=DEBUG
formatter=formatter
args=(sys.stderr,)

[formatter_formatter]
format=%(asctime)s %(name)-12s %(levelname)-8s %(message)s
\end{lstlisting}

\section{Worked Example: Configuring a Two-Backend Workflow}

To make the abstract template above concrete, consider configuring a
workflow that mixes a local \texttt{BATCH} preparation task with a
\texttt{SLURM}-scheduled simulation and an \texttt{LLM}-based summarization
step --- a three-backend combination representative of the synthesis
chapter's forecasting-chain pattern (\Cref{ch:synthesis}). Applying the
``delete what is unused'' principle from the previous section, the correctly
trimmed configuration retains only:

\begin{lstlisting}[style=ini]
[dagon_service]
use = False

[batch]
scratch_dir_base=/data/dagon-scratch/
run_base=%Y%m%d-%H%M%S
threads=4

[slurm]
partition=compute

[llm.research]
endpoint=https://api.example.org
api_key_env=DAGON_RESEARCH_LLM_KEY
model=research-chat
\end{lstlisting}

Note what was deliberately omitted relative to the full template: no cloud
provider section (no \texttt{CLOUD} task is used), no \texttt{ftp\_pub} or
\texttt{dagon\_ip} section (no transversal or service-backed task is used),
and a \texttt{scratch\_dir\_base} deliberately moved off \texttt{/tmp/} to
durable storage, per \Cref{ch:getting-started}'s caution about treating
\texttt{/tmp} as durable across a checkpointed, multi-session operational
deployment. The logging sections are also omitted here on the assumption
that this configuration is supplied as an in-memory \texttt{config}
dictionary rather than loaded via \texttt{config\_file} --- recall from
\Cref{ch:configuration}'s discussion of the layered override mechanism that
\texttt{fileConfig()} is invoked only on the file-loading path, so a purely
programmatic configuration has no logging-section obligation the file-based
path would otherwise impose.

\section{Runtime Environment Variables (not in \texttt{dagon.ini})}

\begin{lstlisting}[style=bash]
; SKYCDS staging credentials
export DAGON_SKYCDS_CLIENT_TOKEN=...
export DAGON_SKYCDS_CATALOG_TOKEN=...
export DAGON_SKYCDS_API_TOKEN=...
export DAGON_SKYCDS_IP=...

; Globus transfer
export DAGON_GLOBUS_TRANSFER_TOKEN=...

; LLM provider credential, named by api_key_env above
export DAGON_LLM_EXAMPLE_API_KEY=...

; Web task authentication, named by token_env / username_env / password_env
export SERVICE_TOKEN=...
\end{lstlisting}

This environment-variable list is, collectively, the concrete instantiation
of the security-expectations discipline reproduced in Appendix~B: every one
of these variables exists specifically so that the corresponding
\texttt{*\_env}-suffixed configuration key (or, in SKYCDS's and Globus's
case, a documented fixed variable name) can name a secret \emph{indirectly},
without that secret ever appearing as a literal string anywhere in
versioned configuration, workflow source, or serialized workflow JSON
(recall \Cref{app:reference}'s note that \texttt{as\_json()} can serialize a
full workflow, including its \texttt{LLM} task provider references --- a
further, concrete reason the credential itself must never be embedded
directly in anything that serialization might capture).

\section{Quick-Reference Table}

\begin{center}
\begin{tabular}{@{}p{0.24\linewidth}p{0.70\linewidth}@{}}
\toprule
\textbf{Section} & \textbf{Purpose} \\
\midrule
\texttt{[dagon\_service]} & Optional remote registration/status service;
  \texttt{use=False} for local-only operation. \\
\texttt{[ftp\_pub]} & FTP host for transversal workflow data download. \\
\texttt{[dagon\_ip]} & Service/node addressing; not required locally. \\
\texttt{[batch]} & Scratch directory base, run grouping, staging thread
  count. \\
\texttt{[slurm]} & Default Slurm partition for staging command generation. \\
\texttt{[llm.<name>]} & One section per OpenAI-compatible LLM provider. \\
\texttt{[ec2]}, \texttt{[digitalocean]}, \texttt{[gce]} & Cloud provider
  credentials for \texttt{CLOUD} tasks (via Apache Libcloud). \\
\texttt{[loggers]}, \texttt{[handlers]}, \texttt{[formatters]}, and
  related sections & Standard Python \texttt{logging.config.fileConfig}
  setup, required when \texttt{Workflow} is constructed without an explicit
  \texttt{config} dictionary. \\
\bottomrule
\end{tabular}
\end{center}

\begin{practicebox}
As a closing reminder consolidating \Cref{ch:configuration}'s scattered
cautions: commit only safe templates (such as \texttt{dagon.ini.sample});
never commit real secrets; record the configuration values that affect
scientific reproducibility (scratch layout, model parameters, container
images, cloud instance types, Slurm resource requests); and keep operational
secrets strictly separate from reproducibility metadata by referencing only
environment-variable \emph{names} in versioned configuration.
\end{practicebox}

\begin{center}
\begin{tabular}{@{}p{0.24\linewidth}p{0.70\linewidth}@{}}
\toprule
\textbf{Section} & \textbf{Purpose} \\
\midrule
\texttt{[dagon\_service]} & Optional remote registration/status service;
  \texttt{use=False} for local-only operation. \\
\texttt{[ftp\_pub]} & FTP host for transversal workflow data download. \\
\texttt{[dagon\_ip]} & Service/node addressing; not required locally. \\
\texttt{[batch]} & Scratch directory base, run grouping, staging thread
  count. \\
\texttt{[slurm]} & Default Slurm partition for staging command generation. \\
\texttt{[llm.<name>]} & One section per OpenAI-compatible LLM provider. \\
\texttt{[ec2]}, \texttt{[digitalocean]}, \texttt{[gce]} & Cloud provider
  credentials for \texttt{CLOUD} tasks (via Apache Libcloud). \\
\texttt{[loggers]}, \texttt{[handlers]}, \texttt{[formatters]}, and
  related sections & Standard Python \texttt{logging.config.fileConfig}
  setup, required when \texttt{Workflow} is constructed without an explicit
  \texttt{config} dictionary. \\
\bottomrule
\end{tabular}
\end{center}

\begin{practicebox}
As a closing reminder consolidating \Cref{ch:configuration}'s scattered
cautions: commit only safe templates (such as \texttt{dagon.ini.sample});
never commit real secrets; record the configuration values that affect
scientific reproducibility (scratch layout, model parameters, container
images, cloud instance types, Slurm resource requests); and keep operational
secrets strictly separate from reproducibility metadata by referencing only
environment-variable \emph{names} in versioned configuration.
\end{practicebox}
